%-------------------------
% Resume - Siddardth Pathipaka
% Version D: Equipment / NPI / Automation Engineer
%-------------------------
\documentclass[letterpaper,10pt]{extarticle}

\usepackage[utf8]{inputenc}
\usepackage[T1]{fontenc}
\usepackage{geometry}
\usepackage{enumitem}
\usepackage{titlesec}
\usepackage{hyperref}
\usepackage{xcolor}
\usepackage{fontspec}
\usepackage{tabularx}

\geometry{
  letterpaper,
  left=0.47in,
  right=0.47in,
  top=0.20in,
  bottom=0.15in,
  noheadfoot
}

\setmainfont{Carlito}[
  BoldFont=Carlito Bold,
  ItalicFont=Carlito Italic,
  BoldItalicFont=Carlito Bold Italic
]

\pagestyle{empty}

\hypersetup{
  colorlinks=true,
  linkcolor={rgb:red,5;green,99;blue,193},
  urlcolor={rgb:red,5;green,99;blue,193},
}
\definecolor{linkblue}{RGB}{5,99,193}

\titleformat{\section}{
  \bfseries\fontsize{11pt}{12pt}\selectfont
}{}{0em}{}[\vspace{-8pt}\rule{\textwidth}{0.5pt}]
\titlespacing*{\section}{0pt}{4pt}{1pt}

\setlist[itemize]{
  leftmargin=0.25in,
  labelsep=0.06in,
  itemsep=0pt,
  parsep=0pt,
  topsep=0pt,
  label=\textbullet
}

\newcommand{\expentry}[2]{%
  \vspace{2pt}%
  \noindent\begin{tabularx}{\textwidth}{@{}X r@{}}
    {\fontsize{11pt}{13pt}\selectfont\textbf{#1}} & {\fontsize{11pt}{13pt}\selectfont\textbf{#2}}
  \end{tabularx}%
  \vspace{-1pt}%
}

\newcommand{\eduentry}[2]{%
  \vspace{2pt}%
  \noindent\begin{tabularx}{\textwidth}{@{}X r@{}}
    {\fontsize{11pt}{13pt}\selectfont\textbf{#1}} & {\fontsize{11pt}{13pt}\selectfont\textbf{#2}}
  \end{tabularx}%
}

\newcommand{\skillline}[2]{%
  \noindent{\fontsize{10pt}{11pt}\selectfont\textbf{#1} #2}\par\vspace{1pt}%
}

\begin{document}

%----------NAME----------
\begin{center}
  {\fontsize{14pt}{16pt}\selectfont\textbf{Siddardth Pathipaka}}\\[2pt]
  {\fontsize{10pt}{12pt}\selectfont\textbf{siddardth1524@gmail.com \textbar\ +1(217) 255-0104 \textbar\ %
  \href{https://github.com/siddardth7}{\color{linkblue}\underline{GitHub}} \textbar\ %
  \href{https://linkedin.com/in/siddardth-pathipaka}{\color{linkblue}\underline{LinkedIn}} \textbar\ %
  \href{https://siddardth7.github.io/portfolio/}{\color{linkblue}\underline{Portfolio}}}}
\end{center}

\vspace{-12pt}

%----------SUMMARY----------
\section{SUMMARY}
{\fontsize{10pt}{11pt}\selectfont
%%SUMMARY%%
}

%----------SKILLS----------
\section{SKILLS}
\vspace{1pt}
%%SKILLS_BLOCK_START%%
\skillline{Equipment \& Tooling:}{Fixture Design, Tooling Qualification, Autoclave \& Cure Equipment, CNC Machining, EDM, HPDC Tooling, GD\&T, CMM Inspection, SolidWorks, AutoCAD}
\skillline{NPI \& Validation:}{FMEA (PFMEA), First Article Inspection, Process Validation, DOE, SPC, DFM awareness, Build Sequencing}
\skillline{Simulation \& Analysis:}{ABAQUS, ANSYS Fluent, FEA (Structural \& Thermal), Classical Lamination Theory, HyperWorks}
\skillline{Programming \& Software:}{MATLAB, Python, C/C++}
%%SKILLS_BLOCK_END%%

%----------EXPERIENCE----------
\section{EXPERIENCE}

% SAMPE -- lead card for NPI/equipment
\expentry{Manufacturing Lead \textbar\ SAMPE Composite Fuselage Program, UIUC}{Jan 2025 \textbar\ May 2025}
\begin{itemize}
  \item Planned and executed a full \textbf{new-part introduction (NPI) sequence} for a \textbf{24-inch composite fuselage} -- from layup tooling setup and process development through \textbf{first-article acceptance} at under \textbf{2\% void content}.
  \item Developed \textbf{PFMEA} for the cure process, identifying vacuum bag integrity as highest-risk failure mode \textbf{(RPN=60)} and configuring automated leak detection -- achieving \textbf{zero process deviations} during production runs.
  \item Performed \textbf{FEA-driven stacking sequence optimization} using ABAQUS and Classical Lamination Theory, reducing structural deflection by \textbf{38\%} while sustaining \textbf{150\% design load (1,000 lbf)}.
  \item Redesigned layup tooling workflow to adopt a \textbf{continuous prepreg blanket} approach, eliminating delamination risk and enabling \textbf{repeatable, scalable} production process.
\end{itemize}

% EQIC -- strong equipment card
\expentry{Manufacturing Intern \textbar\ EQIC Dies \& Moulds Engineers Pvt. Ltd.}{Jun 2022 \textbar\ Aug 2022}
\begin{itemize}
  \item Supported design-to-production build of \textbf{precision die tooling} across a \textbf{12-stage manufacturing sequence} -- CNC machining, EDM, wire EDM, vacuum heat treatment, and final assembly -- tracking dimensional compliance at each stage through to commissioning readiness.
  \item Performed dimensional validation of cavity and core components to \textbf{±0.02 mm} against \textbf{GD\&T specifications} and verified \textbf{parting surface alignment ({>}80\% contact)} on \textbf{800-bar HPDC tooling}, ensuring casting quality at commissioning.
\end{itemize}

% Tata Boeing
\expentry{Manufacturing \& Quality Intern \textbar\ Tata Boeing Aerospace}{Oct 2022 \textbar\ Mar 2023}
\begin{itemize}
  \item Developed and documented a \textbf{CMM inspection process} for \textbf{450+} flight-critical components to \textbf{0.02 mm accuracy}, establishing a repeatable measurement baseline for GE and Boeing program hardware.
  \item Applied \textbf{SPC} to identify \textbf{CNC tool wear} as root cause of recurring position tolerance failures, driving process corrections that reduced defect rate from \textbf{15\% to under 3\%}.
  \item Conducted \textbf{FMEA} on supplier nonconformance, enabling an engineering-approved disposition that prevented \textbf{\textasciitilde\$3,000} in potential scrap.
  \item Drove cross-functional \textbf{8D root cause analysis} for GE engine MRB dispositions, cutting nonconformance resolution cycle time by \textbf{22\%}.
\end{itemize}

% Graduate Research
\expentry{Graduate Research Assistant \textbar\ Beckman Institute, UIUC}{May 2024 \textbar\ Dec 2024}
\begin{itemize}
  \item Developed a novel \textbf{out-of-autoclave cure process} using frontal polymerization that compresses composite processing time from \textbf{8+ hours to under 5 minutes}, demonstrating a scalable alternative to autoclave-constrained production cycles.
  \item Validated process models against experimental measurements within \textbf{10\% of cure front propagation rate} and \textbf{3°C} of peak temperatures, reducing physical validation iterations by \textbf{94\%}.
\end{itemize}

%----------PROJECTS----------
\section{PROJECTS}

\expentry{Brakes \& Suspension Lead \textbar\ Electric Solar Vehicle Championship (ESVC)}{Oct 2019 \textbar\ Apr 2023}
\begin{itemize}
  \item Designed and manufactured brake and suspension hardware for a \textbf{30-member} cross-functional team using manual machining and welding, achieving \textbf{zero mechanical failures} and \textbf{7th place nationally}.
  \item Validated hardware performance with \textbf{180.43/200 on Cross Pad Test}, confirming designed brake and suspension geometry met handling targets.
\end{itemize}

\expentry{Structures \& Manufacturing Lead \textbar\ SAE Aero Design Challenge}{Dec 2021 \textbar\ Sep 2022}
\begin{itemize}
  \item Executed full \textbf{build sequencing plan} for a \textbf{92-inch tapered balsa wing}, managing assembly tooling, reinforcement integration, and \textbf{CG balancing at 25\% MAC} across multi-phase production.
\end{itemize}

\expentry{Research Assistant \textbar\ VNR VJIET Composites Lab, Hyderabad}{Aug 2022 \textbar\ Apr 2023}
\begin{itemize}
  \item Fabricated and tested \textbf{24 composite material formulations} via powder metallurgy, characterizing samples using \textbf{SEM/EDS} and \textbf{Vickers hardness (ASTM E384)} to validate a \textbf{27\% hardness improvement}.
\end{itemize}

%----------EDUCATION----------
\section{EDUCATION}

\eduentry{M.S., Aerospace Engineering \textbar\ University of Illinois at Urbana-Champaign}{Jan 2024 \textbar\ Dec 2025}


\eduentry{B.Tech, Mechanical Engineering \textbar\ VNR VJIET}{Aug 2019 \textbar\ May 2023}

\end{document}